%% mathstc-arbre-probas — arbre pondéré à deux niveaux (deux fournisseurs A et B, pièce
%% défectueuse D). Sur les branches du second niveau on lit des probabilités
%% CONDITIONNELLES : P_A(D) = 0,03. La probabilité d'un chemin est le PRODUIT de ses
%% branches ; P(D) s'obtient en ajoutant les chemins qui aboutissent à D.
\documentclass[tikz,border=4pt]{standalone}
\usepackage{liaisons}
\begin{document}
\begin{tikzpicture}[x=1cm,y=1cm,
  noeud/.style={rounded corners=3pt, draw=black!45, line width=0.7pt, fill=black!8,
                minimum width=0.85cm, minimum height=0.65cm, inner sep=3pt, font=\small},
  branche/.style={line width=0.6pt},
  chemin/.style={solideA, line width=1.8pt},
  poids/.style={font=\scriptsize, inner sep=2pt, fill=white},
  note/.style={font=\scriptsize, text=black!60, inner sep=1.5pt},
  gris/.style={black!55, line width=0.6pt}]

%% ================= les nœuds ==============================================
\fill (0,0) circle (2pt);
\node[noeud] (A)  at (1.55, 1.4) {$A$};
\node[noeud] (B)  at (1.55,-1.4) {$B$};
\node[noeud] (AD) at (3.45, 2.2) {$D$};
\node[noeud] (AN) at (3.45, 0.7) {$\overline{D}$};
\node[noeud] (BD) at (3.45,-0.7) {$D$};
\node[noeud] (BN) at (3.45,-2.2) {$\overline{D}$};

%% ================= premier niveau =========================================
\draw[chemin]  (0.1,0.06) -- (A.west);
\draw[branche] (0.1,-0.06) -- (B.west);
\node[poids, text=solideA] at (0.72,0.95) {$0{,}60$};
\node[poids] at (0.72,-0.95) {$0{,}40$};

%% ================= second niveau ==========================================
\draw[chemin]  (A) -- (AD);
\draw[branche] (A) -- (AN);
\draw[branche] (B) -- (BD);
\draw[branche] (B) -- (BN);
\node[poids, text=solideA] at (2.5,2.05) {$0{,}03$};
\node[poids] at (2.5,0.85) {$0{,}97$};
\node[poids] at (2.5,-0.85) {$0{,}05$};
\node[poids] at (2.5,-2.05) {$0{,}95$};

%% ================= probabilité de chaque chemin ===========================
\node[font=\scriptsize, text=solideA, anchor=west] at (4.0, 2.2) {$0{,}018$};
\node[font=\scriptsize, anchor=west] at (4.0, 0.7) {$0{,}582$};
\node[font=\scriptsize, anchor=west] at (4.0,-0.7) {$0{,}020$};
\node[font=\scriptsize, anchor=west] at (4.0,-2.2) {$0{,}380$};

%% ================= regroupement des deux chemins menant à D ===============
\draw[gris] (4.78,2.2) -- (4.92,2.2) -- (4.92,-0.7) -- (4.78,-0.7);
\draw[gris] (4.92,0.75) -- (5.06,0.75);
\node[note, anchor=west, align=left, text width=2.6cm] at (5.08,0.75)
     {$P(D)=0{,}018+0{,}020$\\$\phantom{P(D)}=0{,}038$};

%% ================= somme des chemins ======================================
\node[note, anchor=north] at (2.3,-2.85) {somme des quatre chemins $=1$};
\end{tikzpicture}
\end{document}
